\section*{Conclusion}
In this paper, we have presented an algorithm able to solve NLS problems subject to constraints of general form. Our method fully exploits the structure of the problem through the AL reformulation and a structured Hessian approximation. 
Although our implementation exhibits promising performance, its limitation for large-scale applications lies in the Hessian approximation: the dense storage of the structured second order correction and the associated computations become prohibitive as the number of variables grows. A limited-memory variant of the structured update, following~\cite{byrd-etal:1994}, would be a natural extension to address this bottleneck. To tackle problems with a large number of residuals, another avenue of research could be to combine the AL framework with an adaptive sampling strategy applied to the residuals, lowering the dimensions of the inner minimization subproblems. This would extend to the constrained case techniques currently investigated for unconstrained problems~\cite{bellavia-etal:2026}.

\vspace{1em}\noindent\textbf{Acknowledgments:} The authors would like to thank Natural Sciences and Engineering Research Council of Canada (NSERC) and the Institute for Data Valorization (IVADO) in addition to Mitacs and  Hydro-Québec for their financial support through grants IT40867 and IT38322.
